\begin{blocksection}
\question \textbf{Tuples.} What would Python display? Run these statements in order. If a statement
causes an error, identify the error and continue with the next statement.
For \lstinline{(3,)} and \lstinline{(3)}, also name the type.

\textbf{Creating tuples}
\begin{lstlisting}
>>> t: tuple[int, str, float, bool, None] = (
...     3, "csm", 2.5, True, None
... )
>>> t
\end{lstlisting}
\begin{solution}[.3in]
\begin{lstlisting}
(3, 'csm', 2.5, True, None)
\end{lstlisting}
\end{solution}
\end{blocksection}

\begin{blocksection}
\begin{lstlisting}
>>> 4, 5, 6, 7
\end{lstlisting}
\begin{solution}[.3in]
\begin{lstlisting}
(4, 5, 6, 7)
\end{lstlisting}
\end{solution}
\end{blocksection}

\begin{blocksection}
\begin{lstlisting}
>>> (3,)
\end{lstlisting}
\begin{solution}[.3in]
\begin{lstlisting}
(3,)
\end{lstlisting}
Type: \lstinline{tuple}.
\end{solution}
\end{blocksection}

\begin{blocksection}
\begin{lstlisting}
>>> (3)
\end{lstlisting}
\begin{solution}[.3in]
\begin{lstlisting}
3
\end{lstlisting}
Type: \lstinline{int}.
\end{solution}
\end{blocksection}

\begin{blocksection}
\begin{lstlisting}
>>> ()
>>> len(())
\end{lstlisting}
\begin{solution}[.45in]
\begin{lstlisting}
()
0
\end{lstlisting}
\end{solution}
\end{blocksection}

\begin{blocksection}
\textbf{Slicing and converting}
\begin{lstlisting}
>>> t[1:4]
\end{lstlisting}
\begin{solution}[.3in]
\begin{lstlisting}
('csm', 2.5, True)
\end{lstlisting}
\end{solution}
\end{blocksection}

\begin{blocksection}
\begin{lstlisting}
>>> t[::-1]
\end{lstlisting}
\begin{solution}[.3in]
\begin{lstlisting}
(None, True, 2.5, 'csm', 3)
\end{lstlisting}
\end{solution}
\end{blocksection}

\begin{blocksection}
\begin{lstlisting}
>>> list(t)
\end{lstlisting}
\begin{solution}[.3in]
\begin{lstlisting}
[3, 'csm', 2.5, True, None]
\end{lstlisting}
\end{solution}
\end{blocksection}

\begin{blocksection}
\begin{lstlisting}
>>> list("csm")
>>> list(())
\end{lstlisting}
\begin{solution}[.45in]
\begin{lstlisting}
['c', 's', 'm']
[]
\end{lstlisting}
\end{solution}
\end{blocksection}

\begin{blocksection}
\textbf{Combining tuples}
\begin{lstlisting}
>>> (4, 5) + (6, 7)
\end{lstlisting}
\begin{solution}[.3in]
\begin{lstlisting}
(4, 5, 6, 7)
\end{lstlisting}
\end{solution}
\end{blocksection}

\begin{blocksection}
\begin{lstlisting}
>>> (3,) * 3
>>> (3) * 3
\end{lstlisting}
\begin{solution}[.45in]
\begin{lstlisting}
(3, 3, 3)
9
\end{lstlisting}
\end{solution}
\end{blocksection}

% Deferred until mutation is introduced; uncomment to restore this prompt.
% \begin{blocksection}
% \begin{lstlisting}
% >>> t[0] = 4
% \end{lstlisting}
% \begin{solution}[.3in]
% \begin{lstlisting}
% TypeError
% \end{lstlisting}
% \end{solution}
% \end{blocksection}

\begin{blocksection}
\textbf{Dictionary keys}
\begin{lstlisting}
>>> locations = {(4, 5): "CSM"}
>>> locations[(4, 5)]
\end{lstlisting}
\begin{solution}[.3in]
\begin{lstlisting}
'CSM'
\end{lstlisting}
\end{solution}
\end{blocksection}

\begin{blocksection}
\begin{lstlisting}
>>> locations[[4, 5]]
\end{lstlisting}
\begin{solution}[.3in]
\begin{lstlisting}
TypeError
\end{lstlisting}
\end{solution}
\end{blocksection}

\begin{blocksection}
\begin{lstlisting}
>>> locations[(4, [5])]
\end{lstlisting}
\begin{solution}[.3in]
\begin{lstlisting}
TypeError
\end{lstlisting}
\end{solution}
\end{blocksection}

% Deferred teaching point: tuple item assignment raises TypeError.
\begin{questionmeta}
    Suggested Time: 8--10 min; Difficulty: Easy--Medium
    \nopagebreak[4]
\begin{itemize}
        \item Ask students to read the annotation on \lstinline{t}: it describes a five-element tuple, with the type of each position listed in order. The annotation does not convert values or enforce their types at runtime.
        \item A comma creates a tuple: parentheses usually group expressions.
        The empty tuple is written \lstinline{()}; one item needs a trailing comma.
        \item Tuples can hold different data types. Slicing returns a tuple;
        \lstinline{list} builds a list from the elements of an iterable.
        \item Tuple addition concatenates and multiplication repeats elements;
        neither operation changes the original tuple. Briefly describe tuples as immutable sequences; focus here on reading and combining values.
        \item Dictionary keys must be hashable. A tuple can be a key only when
        all its elements are hashable. Lists are unhashable, including a list
        nested inside a tuple. Use the last example as a follow-up challenge.
    \end{itemize}
\end{questionmeta}
